% 2012 年江西卷文科第 19 题：梯形 ABCD（图 1）与折叠后的多面体 CDEFG（图 2）。
% 来源 PDF：普通高考/2012/2012江西文.pdf，第 2 页（印面 750）第 2 栏。
% 图 1：AB=12，AE=3，EF=5，FB=4，DE=CF=4，AD=5，BC=4√2，CD=EF=5。
% 图 2：A、B 折起重合于 G，G 在 EF 的垂直平面内（EG=3，FG=4，DG=5，CG=4√2）。
\begin{tikzpicture}[scale=0.62, semithick, font=\small]
  % ---------- 图 1 ----------
  \begin{scope}
    \coordinate (A) at (0,0);
    \coordinate (E) at (3,0);
    \coordinate (F) at (8,0);
    \coordinate (B) at (12,0);
    \coordinate (D) at (3,4);
    \coordinate (C) at (8,4);
    \draw (A) -- (B) -- (C) -- (D) -- cycle;
    \draw (D) -- (E);
    \draw (C) -- (F);
    \node[below left] at (A) {$A$};
    \node[below] at (E) {$E$};
    \node[below] at (F) {$F$};
    \node[below right] at (B) {$B$};
    \node[above left] at (D) {$D$};
    \node[above right] at (C) {$C$};
    \node[below] at (6,-0.8) {图 1};
  \end{scope}
  % ---------- 图 2 ----------
  \begin{scope}[shift={(13.5,0.6)}]
    \coordinate (E2) at (0,0);
    \coordinate (F2) at (5,0);
    \coordinate (D2) at (0,4);
    \coordinate (C2) at (5,4);
    \coordinate (G2) at (1.8,-0.5);
    \draw (E2) -- (F2) -- (C2) -- (D2) -- cycle;
    \draw (D2) -- (E2);
    \draw (C2) -- (F2);
    \draw (G2) -- (E2);
    \draw (G2) -- (F2);
    \draw (G2) -- (D2);
    \draw (G2) -- (C2);
    \node[left] at (E2) {$E$};
    \node[right] at (F2) {$F$};
    \node[left] at (D2) {$D$};
    \node[right] at (C2) {$C$};
    \node[below] at (G2) {$G$};
    \node[below] at (2.5,-2.0) {图 2};
  \end{scope}
\end{tikzpicture}
